% tA_source_facts.tex: generated by analysis/paper_tables.py from analysis/paper_numbers.json; do not edit by hand.
% One tabular plus an optional note; needs booktabs and array. The caption is written in the manuscript.
{\small
\begin{tabular}{>{\raggedright\arraybackslash}p{7.3cm}lll}
\toprule
Fact & gpt-oss:20b & qwen3:8b & qwen3:30b-a3b \\
\midrule
\multicolumn{4}{l}{\textit{Testbed (same for all models)}} \\
Sources per case: resume items / work sample / notes & \multicolumn{3}{l}{5 / 1 / 2 (1152 sources in all)} \\
Gold bait sources that are notes & \multicolumn{3}{l}{144/144} \\
Risk-case sources containing a trap signature that are notes & \multicolumn{3}{l}{126/126} \\
Risk-case notes that are gold bait sources & \multicolumn{3}{l}{144/144} \\
Trap-signature occurrences: in notes only / in a resume item or work sample / in no source & \multicolumn{3}{l}{350 / 0 / 116 (of 466)} \\
Control valid-claim signatures found only in notes & \multicolumn{3}{l}{3/144} \\
Classed subjective or hearsay by the validator's source lexicon: gold baits; control notes; resume items; work samples & \multicolumn{3}{l}{58/144; 22/144; 0/720; 0/144} \\
\midrule
\multicolumn{4}{l}{\textit{C2 supported-side items containing a trap signature (trap items)}} \\
Trap items citing a note & 15/18 & 45/55 & 5/5 \\
Trap items without a source ID & 2/18 & 0/55 & 0/5 \\
Trap items FAILed by any rule & 12/18 & 47/55 & 4/5 \\
\quad of these, FAILed by \texttt{subjective\_claim\_boundary} & 10/12 & 32/47 & 4/4 \\
\quad of these, FAILed by that rule only & 10/12 & 25/47 & 2/4 \\
\bottomrule
\end{tabular}}
\par\smallskip{\footnotesize\raggedright\textit{Note.} Source types come from the case metadata (resume item, work sample, note). Gold bait sources: the gold \texttt{baitSourceIds} of the risk cases; sources containing a trap signature: case-insensitive substring match of a gold trap signature in the source text. Every validator FAIL of \texttt{subjective\_claim\_boundary} on a C2 trap item is the code \texttt{subjective\_source\_cited\_as\_support}. Item level as in the validator table: each FAIL finding is assigned to the deepest supported-side item that contains its path.\par}
